\par\bigskip\noindent\textbf{K101 supporting proof: The complete fixed-locus and curve-clutch ledger}\par\nobreak
This is a working continuation of EI, EIR, IA and FO. The source and
line morphisms are those already written there. The finite local
characteristic calculations below retain their own boundary scope.
The base-direction normal classes and the complete relative
selected-sphere cap integral are still being computed. No source
return is assigned in this note.

\par

\medskip\noindent\textbf{FL.1. Exact extraction from the actual smooth charts}

\par\nobreak

For an ordered unimodular cone matrix N, the physical vector is
\(\delta=(0,0,0,1)\) in the ordinary quotient coordinates
\((\chi^4,\chi\eta,v,V_{\rm phys})\). Its four actual chart
weights are \(N^{-1}\delta\). Set to zero precisely the
coordinates with nonzero weights. This is the fixed chart,
and its smallest cone face is generated by the corresponding
columns of N. Two charts with the same face give two charts of
that fixed component. A recorded period word identifies the
faces only by its actual integral ray map and full coefficient
map. These statements follow directly from the diagonal chart
action, not from the dimension of a cone alone.

For a source and target chart, the complete period map has
exponent matrix \(N_t^{-1}T N_s\). Its unit vectors are
\(N_t^{-1}(0,1,-1,0)^t\) for the entire source unit \(c_b\)
and \(N_t^{-1}(0,0,0,1)^t\) for
\(\mathcal C=\lambda H\mathcal B\). On a fixed overlap,
every normal target row has one normal source coordinate to
the first power, with the remaining factors in invertible
tangent coordinates and the specified units. Restricting
these factors to the fixed component is the actual normal
bundle transition. Derivatives of the full units are retained
in the total chart map; terms multiplying a vanishing normal
coordinate have exactly this stated first derivative.

The executable record
\path{code/part3_fourfold/e6_fixed_face_orbits.py} recomputes every
weight and every normal row, and saves the complete parent,
period and unit maps. For the 238 ordinary O/T quotient
charts it gives the following local list, for one O or T
source neighbourhood:

\begin{longtable}[]{@{}
  >{\raggedright\arraybackslash}p{(\columnwidth - 4\tabcolsep) * \real{0.3000}}
  >{\raggedleft\arraybackslash}p{(\columnwidth - 4\tabcolsep) * \real{0.4000}}
  >{\raggedright\arraybackslash}p{(\columnwidth - 4\tabcolsep) * \real{0.3000}}@{}}
\toprule\noalign{}
\begin{minipage}[b]{\linewidth}\raggedright
Component charts
\end{minipage} & \begin{minipage}[b]{\linewidth}\raggedleft
Number
\end{minipage} & \begin{minipage}[b]{\linewidth}\raggedright
Retained boundary
\end{minipage} \\
\midrule\noalign{}
\endhead
\bottomrule\noalign{}
\endlastfoot
Interior fixed points & 106 & None in this local fan \\
Interior fixed curves & 18 & Complete two-chart stars \\
Fixed curves reaching the \(\chi\)-unit overlap & 10 & The 38-cone interface \\
Fixed curves reaching the v-unit source chart & 12 & The other source blowup chart \\
Strict-source surface patches & 4 & \(\chi\)-unit and \(\eta\)-unit \\
Exceptional-source surface patches & 4 & \(\chi\)-unit and v-unit \\
Refinement surface patches & 10 & \(\chi\)-unit \\
\end{longtable}

The finite boundary tags use actual nonnegative cover
valuations: \(4\operatorname{ord}\chi=n_0\),
\(4\operatorname{ord}\eta=4n_1-n_0\), and
\(\operatorname{ord}v=n_2\). Thus a tag specifies a
literal unit-coordinate overlap, not an inferred compactness
property. Both O and T copies retain their own source charts.
The period B interchanges these two copies while increasing
the chain index by four. The complete source map is still
the odd-b chord, including its projective charts and units.

The same extraction on FO\textquotesingle s 39 selected charts identifies
all 18 surface patches and all 10 interface curves with
the 38-cone common interface. Its one additional surface
face is
\[
 \langle Q_0,Q_1\rangle,
 \qquad Q_0=(0,0,1,0),\quad Q_1=(0,0,1,1)
\tag{FL.1}
\]
in \((\eta,v,q,V_\theta)\). It is the moving selected
fixed surface, with \(q=0\). Its two physical normal
weights are \((-1,+1)\); its tangent coordinates are
\(\eta,v\), with \(\sigma=\eta v\). No boundary of
this surface is counted as an interior compact component.

The actual interface is the one in FO.6:
\[
 U_{\rm selected}=q^{-j}\psi\,(\eta v)^{-4j}V_{\rm phys},
 \quad q=\kappa H,\quad
 \beta=\mathcal B\frac{\psi(T_bp)}{\psi(p)}.
\tag{FL.2}
\]
Hence the selected multiplier \(\lambda\beta/\kappa\)
and the ordinary multiplier \(\lambda H\mathcal B\)
give the same B square, including the complete \(q^j\)
and \(\psi\) factors. The monomial parent maps of the
39- and 238-cone refinements realize this square on all
their common faces.

\par

\medskip\noindent\textbf{FL.2. The four strict-source patches and their actual source target}

\par\nobreak

Write
\[
 C_j=(0,0,1,j),\qquad A_j=(0,1,0,j),\qquad0\le j\le4.
\]
The faces \(\langle C_j,C_{j+1}\rangle\) are precisely
the four v-zero strict-source patches. To see the source
map directly, use the second factor chart
\[
 v=st,\qquad X_j=\eta s,\qquad Y_j=t,
 \qquad\sigma=\eta st.
\tag{FL.3}
\]
Its fixed surface is s=t=0, with \((\chi,\eta)\)
as source coordinates. On it the source equation is
\[
 \bar Y^2=\bar X^3-2\chi^2\bar X^2+\bar X.
\]
The actual fourth-root quotient map is
\[
 t_e=\chi^{-4},\qquad x_e=\chi^{-2}\bar X,
 \qquad y_e=\chi^{-3}\bar Y,
\tag{FL.4}
\]
and it satisfies
\(y_e^2=x_e^3-2x_e^2+t_ex_e\) identically. At the
infinity chart the same morphism is polynomial:
\[
 q_\infty=\chi^4,\qquad x_0=\chi^6\bar X,
 \qquad y_0=\chi^9\bar Y,
\]
\[
 y_0^2=x_0^3-2q_\infty^2x_0^2+q_\infty^3x_0.
\tag{FL.5}
\]
Each coordinate is invariant under the actual fourth-root
action. These maps give the concrete Weierstrass target
of every strict-source surface patch. The finite-native
I2 resolution is joined by IA.5-\/-13, including t zero.
The selected fibre is joined by FL.2 and the six/39
charts. The quotient modifications at \(\chi=0\)
retain their invariant polynomial maps and original
volume characters. Thus no identification of a modified
surface with an unmodified minimal surface is needed to
specify its source projection or normal maps.

On the native finite side, before the source graph
blowup, write the line coordinates as
\(A_jB_j=\epsilon\). Near an O or T section the
strict-source graph chart has
\[
 A_j=\eta_e s,\qquad B_j=t,\qquad
 \epsilon=\eta_e st.
\tag{FL.6}
\]
Consequently the positive normal frame is changed by
the actual section factor \(\eta_e\), while the
negative frame is unchanged. The nonzero point-line
frame is the actual \(\sigma_a/r_a\), not the
meromorphic section by itself. Formula IA.12 carries
that frame to the complementary source chart, and
IA.13 supplies every \(\chi^j\) factor. These are
the base-direction normal transitions still to be
combined with the finite fibre charts in the global
normal-class calculation.

\par

\medskip\noindent\textbf{FL.3. The exceptional and refinement surface fibres}

\par\nobreak

The faces \(\langle A_j,A_{j+1}\rangle\) are the
exceptional-source patches. In the first factor chart
\[
 \eta=st,\qquad X_j=s,\qquad Y_j=vt,
 \qquad\sigma=vst.
\tag{FL.7}
\]
The fixed surface has coordinates \((\chi,v)\).
The other source blowup chart is \(\eta=\sigma p\)
with \(p=1/v\). Its line-chain coordinates are
\[
 A=s,\qquad B=vt,
 \qquad AB=\sigma.
\tag{FL.8}
\]
This is the full raw-parent normal arrow along the projective
exceptional source fibre, with inverse \(s=A\),
\(t=pB\). In the ordered negative/positive physical
normal lines its raw-parent degrees on that fibre are \((-1,0)\).
Indeed the negative coefficient changes by v, which
is the standard \(\mathcal O(-1)\) clutch in these
specified coordinates, and the positive coefficient
is unchanged. FL.8 records the additional normal factors
in the actual common refinement.

The point-line factor through the whole exceptional
fibre is retained: in the p-chart
\(r_a=(p-c_a)/p\), and the nonzero line frame is
\(\sigma_a/r_a\). Its zero and pole are cancelled
by the actual point and O Cartier sections. At p
infinity it agrees with the canonical frame. Therefore
FL.8 is stated in that full frame, including the
points p zero and p equals \(c_a\), rather than in
a punctured meromorphic coefficient chart. The full
odd-b reset sends p to p minus b and interchanges
O and T; the coefficient remains the exact Hom unit
in that frame. Period powers and the finite quotient
are retained at the two ends of the base direction.

Each of the ten remaining surface faces lies over
\(\eta=v=0\), and its \(\chi\)-unit fibre is a
complete toric \(\mathbb P^1\). The two adjacent
three-cones give its whole normal clutch. In the
literal sorted ray order, with weights \((-1,+1)\),
the fibre degrees are:

\begin{longtable}[]{@{}ll@{}}
\toprule\noalign{}
Face in \((\eta,v,V)\) & Normal degrees \\
\midrule\noalign{}
\endhead
\bottomrule\noalign{}
\endlastfoot
\((1,1,2),(1,1,3)\) & \((-1,-1)\) \\
\((1,1,3),(1,1,4)\) & \((-2,-1)\) \\
\((1,1,4),(1,1,5)\) & \((-1,-2)\) \\
\((1,1,5),(1,1,6)\) & \((-1,-1)\) \\
\((1,2,3),(1,2,4)\) & \((-1,-1)\) \\
\((1,2,4),(1,2,5)\) & \((0,-1)\) \\
\((1,2,5),(1,2,6)\) & \((-1,0)\) \\
\((2,1,6),(2,1,7)\) & \((0,-1)\) \\
\((2,1,7),(2,1,8)\) & \((-1,0)\) \\
\((2,1,8),(2,1,9)\) & \((-1,-1)\) \\
\end{longtable}

The certificate \path{e6_selected_infinity_surface_fibres.py}
computes every degree from the two full inverse cone
matrices with two independent auxiliary cocharacters.
Their tangent weights are opposite and nonzero; the
normal degree is the sum of normal weight divided by
tangent weight at the two endpoints. Equivalently it
is the degree of the explicitly saved monomial line
transition. The two computations agree in each line.
They retain the signed physical normal ordering.

These fibre computations leave the base-direction
normal clutch explicit. In particular they do not
replace a ruled surface by a product, nor set its
base Chern coefficients to zero.

\par

\medskip\noindent\textbf{FL.4. The finite source patches and their local characteristic terms}

\par\nobreak

All 106 interior point faces and all 18 interior
curve faces of the 238-cone patch are contained over
the actual local source origin. The curve stars
have two charts with opposite tangent weights;
the complete normal degree agrees for two generic
auxiliary circles. Units pulled back from the
source origin restrict to their same constant on
each such compact fibre. Their complete normal
derivative is therefore the one recorded by the
monomial-and-unit chart maps in FL.1.

Using SZ.57-\/-59 in its positive-hyperplane convention,
these compact interior components have total local
value \(9/4\) for one O or T patch. The actual
physical line is \(M=U^{-1}\), so both its cubic
and its Pontryagin-linear integrals change sign;
the same local value for M is \(-9/4\). This
calculation is a local part of the boundary ledger.
The O and T copies and their global surface terms
have not been collapsed into a source return.

The additional two-torsion pair is carried by the
root\textquotesingle s explicit fixed-ideal quotient blowup. Its
four surface patches have the literal A1 map
\[
 (\rho,z)\longmapsto(\chi^2,\chi\zeta,\zeta^2)
                  =(\rho,\rho z,\rho z^2).
\tag{FL.9}
\]
The common boundary is
\(|\rho|(1+|z|^2)=r^2\), with the same complex
outward orientation. The odd normal line map is
\(\ell=\chi a\), and its order-two boundary
holonomy is retained. The integral intersection
matrix is \([-2]\). It gives
\(\Delta D^2=-1/2\), where
\(D=c_1(L_+)-c_1(L_-)\). The extending determinant
form and the same outward vector field give
\(\Delta p_1=-3\). Hence its signed surface change
is \(-1/16\), cancelling the isolated-point
\(+1/16\) in the positive-hyperplane convention.
The physical dual negates both, retaining zero.

The complete source mixing, \(\gamma^{-1}B\) and
line maps commute with FL.9. The source/normal
boundary morphisms are the ones in
\path{part3_e6_selected_infinity_mu2_residue_comparison_working.md};
their full independent peer audit is
\path{records/fourfold_audit_e6_selected_infinity_mu2_residue.json}.
Thus this zero correction can be used with the
original orbifold surface patches without losing
their boundary line character. Their remaining
global surface integrals are still retained.

\par

\medskip\noindent\textbf{FL.5. Complete geometric cover and the remaining relative cap data}

\par\nobreak

Use the actual proper source-base graph
\(t_o q_{\infty,e}=\epsilon^4d^{-3}\). The following
overlaps specify the finite cover of its completion:

\begin{longtable}[]{@{}
  >{\raggedright\arraybackslash}p{(\columnwidth - 2\tabcolsep) * \real{0.5000}}
  >{\raggedright\arraybackslash}p{(\columnwidth - 2\tabcolsep) * \real{0.5000}}@{}}
\toprule\noalign{}
\begin{minipage}[b]{\linewidth}\raggedright
Pieces
\end{minipage} & \begin{minipage}[b]{\linewidth}\raggedright
Actual overlap
\end{minipage} \\
\midrule\noalign{}
\endhead
\bottomrule\noalign{}
\endlastfoot
Native finite I2 source and complementary source & IA.5-\/-13; both source graph ideals multiply by the same unit \(\chi\) \\
Ordinary and selected O/T line graphs & FO\textquotesingle s 38-cone common interface, its 39/238 extensions, and FL.2 \\
Native and complementary ordinary chains & \(X_k=\chi^kA_k\), \(Y_k=\chi^{-k-1}B_k\), with unchanged physical coefficient \\
Complementary graph and original ordinary exterior & \(t_o=\sigma^4d^{-3}\), \(x_o=\sigma^2d^{-2}\bar X\), \(y_o=\sigma^3d^{-3}\bar Y\), and the canonical point-section arrow \\
Ordinary exterior and original finite cap & The unchanged K96 sector maps \(\Phi_\beta=A^kB^{j_{\alpha\beta}}\Phi_\alpha\) \\
Fourth-root graph and that cap & \(\gamma\) lifts to \(AB^{j_\gamma}\); four sectors give \(A^4B^{\sum j_\gamma}\) \\
Ordinary finite stabilizers & Both 238 O/T quotient maps and the full two-torsion fixed-ideal blowups, with all source units and determinant characters \\
\end{longtable}

The source graph modifications are proper. The common
factor refinements are proper over their original
factor parents by their full-support projective graph
maps. On the ordinary source regions the actual Hom
coefficient is a nonzero bounded unit times
\(\lambda\sigma^4H\), or its native
\(\lambda\epsilon^4\widehat B\) expression. Hermitian
norms in the actual point-line atlas retain those
bounds on closed source regions. The line-chain
period shifts by four nodes. The selected two-period
fundamental region has the exact positive determinant
and bounded coordinates IA.1-\/-2. Together these give
finite compact fundamental chart regions over the
closed source-base graph, with all adjacent labels
and their actual unit maps retained.

On the original exterior and cap the source parameter
family is the extending K96 family. Their overlaps
are unchanged by modifications supported on the
central graph. On every triple overlap the displayed
maps are the same source/point-line maps on the dense
ordinary open set; their holomorphic extensions in
the written normal charts preserve that equality.
The source and line components of each gluing, and
the entire cap reference cocycle, are therefore fixed
by the ledger. The residual normal-class calculation
must use these maps rather than insert another
annular period.

The only moving fixed surface with an outer parameter
boundary is FL.1 continued along the selected
normalization. Its boundary is the actual fixed
sphere family. To form the relative normal-link
calculation, retain its precise normal clutch to
the chosen reference cap. In the native fixed-edge
Cox chart the map to the original chart is
\[
 (b,a,c)\longmapsto
 (b,\epsilon^4d^{-3}\chi_L^{-1}a,\chi_Lc),
 \qquad\chi_L=(\epsilon/d)^{1-a_{
m line}}.
\tag{FL.10}
\]
Its positive-weight coefficient map is the actual
conjugate of the first normal line together with
the unchanged complex second line. The complete
point/tensor coefficient is EI.20-\/-27, with all
theta and Miller factors; the source normalization
map is EI.8, including its two projective graph
charts. The physical Hopf identification is the
K68 anti-linear map with its original parameter
and inverse-clutch direction. These maps specify
the boundary bundle and its framing. Its common
parameter phase, its relative SU2 factor and the
source projective change must all be carried in
the reference-cap comparison.

The remaining characteristic tasks are now precise:
the base-direction normal classes of the four
strict-source, exceptional-source and refinement
surface continuations, and the full normal clutch
of FL.10 across that selected-sphere reference cap.
The ten interface curves also retain their two
end normal maps under the native/selected chart
comparison. The local compact residues in FL.4
and the finite two-torsion cancellation do not
assign these missing integrals. No global
characteristic sum or source return is asserted.

\par

\medskip\noindent\textbf{FL.6. The full base clutch of the ten interface curves}

\par\nobreak

This section completes the ten curve maps left in FL.5. It leaves
the surface and moving-cap calculations stated there unchanged.
Use the actual native finite source coordinate \(\eta_e=-x_e/y_e\)
at O. At T use \(b_0=y_e/t_e\) in IA.5; its resolved source
equation has a unit derivative in the solved variable a
at the central section, including \(t_e=0\). Write either local
source coordinate as \(\eta_e\) in the following common graph
formula. The source graph is the blowup of
\((\eta_e,\epsilon)\), and its chosen chart has
\(v=\epsilon/\eta_e\). On the entire complementary overlap,
IA.12 gives the exact map
\[
 \eta_e=\chi\eta,\qquad v_e=v,\qquad
 V_e=V_{\rm phys},\qquad t_e=\chi^{-4}.
 \tag{FL.11}
\]
For T the same equality is \(b_0=\chi\bar Y\), with the
identical v ratio. Thus the three native factor coordinates
are literally \((T_1,T_2,T_3)\) in the actual quotient torus
\[
 (T_0,T_1,T_2,T_3)=(\chi^4,\chi\eta,v,V_{\rm phys}).
\]
This map carries the complete Cartier frame: the ratio
\((\tau-\tau(aG))/\tau\) is unchanged because both formal
coordinates multiply by \(\chi\). Away from its two Cartier
divisors the canonical point section supplies the same arrow.
No tensor-frame parameter power is inserted into FL.11.

For one of the ten fixed curves let its three sorted common
rays form the matrix \(F\) in \((\eta,v,V)\), and set
\(A=F^{-1}\). The native normal coordinate characters are
the rows of \(A\) applied to \((T_1,T_2,T_3)\). The complete
four-dimensional exponent matrix is therefore
\[
 E=(0\mid A).
 \tag{FL.12}
\]
The recorded 238-cone chart containing its infinity endpoint
has ordered matrix \(N\). Direct multiplication shows that
each row of \(EN\) is one normal coordinate to exponent one,
times a power \(d_i\) of the sole tangent coordinate. The
first row of \(N\) is the unit vector of this tangent column;
hence that coordinate is exactly \(T_0=q_\infty\), without
an additional normal factor. The actual line arrows are
\[
 n_{e,i}=q_\infty^{d_i}n_{\infty,i}.
 \tag{FL.13}
\]
The inverse uses the reciprocal monomial on the same punctured
overlap. With finite coordinate \(t_e\) and
\(q_\infty=t_e^{-1}\), this coefficient clutch has degree
\(-d_i\). For example the coefficient relation
\(n_e=q_\infty n_\infty\) is the tautological line clutch
from the frames \((1,t_e)\) and \((q_\infty,1)\), and has
degree minus one. This fixes the sign in every row of FL.13.

The finite native charts extend over the whole finite
\(t_e\)-line. At \(t_e=0\) the literal IA.5-\/-10 I2 charts
provide their continuation. At the selected fibre, FO.26
cones \(Q_0\) over this same three-cone. Every row of the
native normal exponent matrix has zero coefficient on that
q tangent ray. The full point/theta comparison is still
\(V_\theta=\psi V_{\rm phys}\), with
\(\psi=f_{P^a}/r_a\). On \(v=0\) and away from the O
section, \(\sigma=0\), \(aG=O\), and \(P^a=1\).
Here the actual marked Tate representative of G is the
holomorphic unit branch limiting to minus one on the selected
nodal normalization. Since G is two-torsion on \(\sigma=0\),
its square is \(Q^h\) for an integer h. Its valuation zero
at \(q=0\) forces h zero. The retained branch therefore
has \(P_G=-1\) on the full selected disc and
\(P_G^a=1\) for both admitted even a, with no source-period
reindex. The actual theta ratio is then exactly \(f_1=1\), and
the point Cartier ratio is \(r_a=1\). The previously
proved holomorphic nonzero extension of \(\psi\) therefore
has value one also at \(\eta=0\). At T the canonical
point-section chart gives \(r_a=1\) directly and the same
identity holds. Thus all the recorded powers of \(\psi\)
in the normal map restrict to one on the curve. This is
an equality of the full unit restricted to the component.
The complete off-component function remains in the chart map.
The derivative terms of that unit multiply a vanishing
normal coordinate, so they do not add another normal clutch.

These maps give an explicit compact curve: its finite
\(t_e\)-chart and its single \(q_\infty\)-chart glue by
\(q_\infty=1/t_e\). Its entire normal bundle is FL.13.
Each face has no nonzero period stabilizer; its O and T
representatives remain the two specified fundamental
copies. The odd-b source translation and its line unit
identify consecutive period blocks by the recorded B map.
They introduce no additional quotient on these interior
faces. The original finite cap is on the large-sigma
domain and does not meet these sigma-zero curves.

The exact rows, in their physical normal order, are:

\begin{longtable}[]{@{}
  >{\raggedright\arraybackslash}p{(\columnwidth - 6\tabcolsep) * \real{0.2500}}
  >{\raggedright\arraybackslash}p{(\columnwidth - 6\tabcolsep) * \real{0.2500}}
  >{\raggedright\arraybackslash}p{(\columnwidth - 6\tabcolsep) * \real{0.2500}}
  >{\raggedright\arraybackslash}p{(\columnwidth - 6\tabcolsep) * \real{0.2500}}@{}}
\toprule\noalign{}
\begin{minipage}[b]{\linewidth}\raggedright
Face IDs
\end{minipage} & \begin{minipage}[b]{\linewidth}\raggedright
Physical weights
\end{minipage} & \begin{minipage}[b]{\linewidth}\raggedright
Normal degrees
\end{minipage} & \begin{minipage}[b]{\linewidth}\raggedright
Positive-u value per curve
\end{minipage} \\
\midrule\noalign{}
\endhead
\bottomrule\noalign{}
\endlastfoot
1, 3, 5, 54, 68 & \((-1,-1,1)\) & \((-1,-1,0)\) & \(-1/4\) \\
13, 27, 35 & \((1,1,-1)\) & \((1,0,-1)\) & \(1/4\) \\
84, 85 & \((1,1,-1)\) & \((0,1,-1)\) & \(1/4\) \\
\end{longtable}

For completeness the curve term of SZ.57, obtained by
expanding the three actual normal factors, is
\[
 R_u(C)=-\frac{1}{24\prod_i w_i}
 \sum_i\left(\frac{4-\sum_jw_j^2}{w_i}+2w_i\right)
            \deg N_i.
 \tag{FL.14}
\]
It gives five negative quarters and five positive quarters.
The sum is zero in each O or T fundamental copy. The
physical line \(M=U^{-1}\) negates each displayed
individual value and retains the same zero sum.

The reproducible certificate
\path{code/part3_fourfold/e6_selected_infinity_curve_clutch.py}
passes 183 exact checks. Its JSON retains every full
matrix \(E,N,EN\), physical weight, selected q row,
point/theta unit exponent, native correspondence and input
hash. The original local component extraction is unchanged.
The surface integrals and the selected reference-cap
normal clutch remain the outstanding global calculation.

\par

\medskip\noindent\textbf{FL.7. Strict-source normal characters on the full overlap}

\par\nobreak

In the native chain and strict graph use
\(A_jB_j=\epsilon\), \(A_j=\eta_e s_e\),
\(B_j=t_e\), and \(\epsilon=\eta_e s_et_e\).
Here \(t_e\) denotes a normal coefficient; the source-base
coordinate is held fixed. In the complementary graph use
\(X_jY_j=\sigma\), \(X_j=\eta s\), \(Y_j=t\),
and \(\sigma=\eta st\). Substituting the actual overlaps
\(\eta_e=\chi\eta\), \(X_j=\chi^jA_j\), and
\(Y_j=\chi^{-j-1}B_j\) gives
\[
 s=\chi^{j+1}s_e,\qquad t=\chi^{-j-1}t_e,
 \qquad 0\le j\le3.
 \tag{FL.15}
\]
Both products and the exact parameter \(\epsilon=\chi\sigma\)
are preserved. The finite coefficient map \(A_j=\eta_e s_e\)
is the actual inclusion of the ideal of the O or T section;
the negative coefficient map is the identity. Thus the
positive normal line retains that section ideal, and FL.15
retains the two opposite fourth-root characters \(j+1\)
and \(-j-1\). The full point frame remains
\(\sigma_a/r_a\), whose point and parameter cancellation
is the one in FL.11. The source transition is still the
whole odd-b chord with its projective continuation.

These exact coefficients are the current input for the
strict-source surface calculation. Combining them with
the remaining exceptional and selected charts to compute
the global normal squares and the reference-cap clutch
is still in progress. No surface or return value is
assigned here.

\par

\medskip\noindent\textbf{FL.8. The actual common-refinement normal factors}

\par\nobreak

The raw-parent arrows in FL.3 and FL.7 are retained above.
The common38 fan has further normal factors along the O/T
sections. The following calculation is their actual parent
morphism and corrects any use of the raw-parent degrees
as degrees on the final fixed surfaces. The prior working
text and certificate are preserved in
\path{records/attempts/e6_selected_infinity_refined_surface_normals/}.

For the strict face \(C_j,C_{j+1}\), the unique third
ray in the common3 fan is respectively
\[
 (1,3,4),\quad(1,2,6),\quad(1,1,6),\quad(1,0,4).
 \tag{FL.16}
\]
Write it \((1,k,l)\), and let a be its tangent
coordinate. The complete source and line monomials are
\[
 \eta=a,\quad v=a^k n_-n_+,\quad
 V=a^l n_-^j n_+^{j+1}.
\]
In the alternating v-factor parent,
\(s_{\rm old}=V/(\eta^{j+1}v^j)\),
\(t_{\rm old}=v/s_{\rm old}\). Literal substitution gives
\[
 s_{\rm old}=a^{3-j}n_+,\qquad
 t_{\rm old}=n_-.
 \tag{FL.17}
\]
Thus the refined positive line gains the displayed
\(3-j\) powers of the section ideal. If \(L\) denotes
the base-root line whose frame satisfies
\(e_\infty=\chi^{-1}e_{\rm fin}\), the exact source
and point-line arrows give the following orbifold
normal extensions before the local quotient resolution:
\[
 N_+=\mathcal I_{O+T}^{\,4-j}\otimes L^j,
 \qquad N_-=L^{-j-1}.
 \tag{FL.18}
\]
Indeed at O or T the ideal frame changes by
\(\eta_e=\chi\eta\) or \(b_0=\chi\bar Y\).
Together with FL.15 and FL.17 the refined normal
coefficients change by \(\chi^4\) and
\(\chi^{-j-1}\). Away from the two section ideals
the actual chain characters are \(\chi^j\) and
\(\chi^{-j-1}\). These are exactly the two frame
descriptions of FL.18. The actual point-line central
trivialization and the full \(\psi\) comparison in
FL.11-\/-13 fix the remaining scalar; no added annular
cochain is used.

For the exceptional face \(A_j,A_{j+1}\), the third
rays are respectively
\[
 (0,1,0),\quad(1,1,2),\quad(2,1,6),\quad(3,1,12).
 \tag{FL.19}
\]
Writing the tangent coordinate as v gives the
literal parent maps
\[
 s_{\rm old}=n_+,\qquad t_{\rm old}=v^j n_-.
\]
The second source graph chart in FL.8 therefore has
\[
 A=n_+,\qquad B=v^{j+1}n_-.
 \tag{FL.20}
\]
Consequently the refined exceptional-fibre normal
degrees in negative/positive physical order are
\((-j-1,0)\). This follows from the specified
\(p=1/v\) coordinate clutch, including its inverse;
the point-section factors remain in their actual
Cartier frames. The earlier pair \((-1,0)\)
is its raw-parent value and is the refined value
at j zero.

The strict surface star in the 238 chart has tangent
rays \((0,1),(1,1),(2,1),(3,1),(4,1)\) under the
actual source projection \((T_0,T_1)\). The normal
degrees of the two physical lines on its three
exceptional curves, ordered negative then positive,
are
\[
\begin{array}{c|ccc}
 j& C_1&C_2&C_3\\ \hline
 0&(1,0)&(0,0)&(0,0)\\
 1&(1,0)&(-1,0)&(1,0)\\
 2&(0,0)&(0,0)&(1,0)\\
 3&(0,0)&(0,0)&(0,0).
\end{array}
 \tag{FL.21}
\]
These degrees are obtained from the two adjacent
full four-cone normal monomials at each curve.
The integral source intersection matrix has diagonal
minus two and adjacent entries one. Solving that
integral map for the actual difference line
\(D=c_1(N_+)-c_1(N_-)\) gives its relative squares
\(-3/4,-1,-3/4,0\), respectively. Its boundary
characters are retained by the full quotient
monomials and FL.17, so these are relative
comparisons to the specified orbifold lines.
The tangent comparison and all global normal
squares are being assembled with the same
boundary arrows; no global return is assigned.

